\begin{scanpage}{0179}{161}
1. By the method of an integral equation for the current on the surface of a good conducting body (the local character of the field in the half-shade region, see [17]);

2. By the method of asymptotic summing of diffraction series (the current on a paraboloid [64, 65]; the propagation of radio waves above the spherical Earth [18, 66]);

3. By the parabolic equation method (the propagation of radio waves above the flat [67] and spherical [68, 69] Earth; the field of a plane electromagnetic wave in the half-shade region for any convex body [70]).

Keller [42], basing his work on the diffraction ray concept, obtained an expression for the field in a deep shadow with diffraction by a convex cylinder with a variable curvature. In the particular cases of an elliptic and a parabolic cylinder, as was shown in the works of Vaynshteyn and Fedorov [71] and of Ivanov [35], Keller's equations agree with the results of the more rigorous mathematical investigation. This allows one to more precisely study (see [71]) the conversion of diffraction rays to ordinary rays and vice versa.

The parabolic equation method described in so-called ray coordinates is a more general approach to diffraction by convex bodies. This method allows one to obtain a general expression for the Green function in the case of a circular cylinder [72, 73]. Evidently this method can subsequently be successfully applied also to other cases, among them three-dimensional diffraction problems.
\end{scanpage}
